\documentclass[12pt]{article}
\usepackage[utf8]{inputenc}
\usepackage[T1]{fontenc}
\usepackage[margin=1in]{geometry}
\usepackage{setspace}
\usepackage{xcolor}
\usepackage[colorlinks=true, urlcolor=blue!60!black, linkcolor=blue!60!black]{hyperref}
\setlength{\parindent}{0em}
\setlength{\parskip}{0.6em}
\renewcommand{\baselinestretch}{1.0}
\begin{document}
\begin{center}
{\LARGE \textit{DeShaney v. Winnebago County Department of Social Services}, 489 U.S. 189 (1989)\par}
\vspace{0.5em}
{\large Case Brief --- Two-Page Edition\par}
\end{center}

\section{Memory Jogger}
The Due Process Clause imposes no general duty on the State to protect people from private violence, even when social workers know a child is being abused (\href{https://supreme.justia.com/cases/federal/us/489/189/}{489 U.S. 189}).

\section{Detailed Case Facts}
Wisconsin's Winnebago County social-services agency received repeated reports that Randy DeShaney was abusing his son Joshua, briefly placed Joshua in the hospital's temporary custody, and then returned him under a voluntary agreement. A caseworker recorded continuing signs of abuse but did nothing. In 1984 Randy beat four-year-old Joshua into a coma, leaving him permanently brain-damaged.

\section{Procedural History}
Joshua and his mother sued under 42 U.S.C. \S~1983; the district court granted summary judgment for the county, and the Seventh Circuit affirmed.

\section{Judicial Votes}
Affirmed, 6--3. Rehnquist wrote for White, Stevens, O'Connor, Scalia, and Kennedy; Brennan, joined by Marshall and Blackmun, dissented, and Blackmun also dissented separately.

\section{Holding}
The Clause limits state action and creates no affirmative right to protection from private harm, except where the State restrains liberty through custody or creates or worsens the danger.

\section{Analysis of Opinions}
Rehnquist read the Clause as protecting ``the people from the State, not to ensure that the State protected them from each other.'' \textit{Estelle} and \textit{Youngberg} impose duties only when the State holds a person in custody against his will; Joshua was in his father's custody, and returning him left him no worse off. The Amendment ``does not transform every tort committed by a state actor into a constitutional violation,'' and removing Joshua early would have invited parental claims. Brennan argued that Wisconsin's system funneled all abuse reports to DSS and then left Joshua trapped. Blackmun decried ``sterile formalism'': ``Poor Joshua!'' The ``played no part'' sentence later spawned the circuits' state-created-danger doctrine.

\section{Examples: Future Applications}
\textbf{Same-Side:} police who respond slowly to calls about a protective-order violation face no due process claim (\textit{Castle Rock}). \textbf{Opposite-Side:} a trooper who strands a child on an unlit interstate at night created the danger. \textbf{Fence-Sitter:} abuse of a child in a state-selected foster home, which footnote 9 left open.

\section{Critique}
Critics, echoing Brennan, say ``inaction'' is a baseline choice when the State gave Randy custody and monopolized rescue. The ruling also sits uneasily with the Fourteenth Amendment's origins in the States' refusal to protect Black citizens from Klan violence.

\section{Key Quotations}
\begin{quote}
``While the State may have been aware of the dangers that Joshua faced in the free world, it played no part in their creation, nor did it do anything to render him any more vulnerable to them.''
\end{quote}
\noindent --- Chief Justice Rehnquist, for the Court, 489 U.S. at 201.

\end{document}
